\chapter{Radicals}\label{ch:rad}

$\sqrt 8 = 2\sqrt 2$ is a fact every student knows and few can say precisely. What is $\sqrt 8$?
Why is $2\sqrt 2$ simpler? What is the general rule, and for cube roots? The chapter answers
these from the fundamental theorem of arithmetic, then meets the one place in the book where the
engine cannot write the textbook's answer --- and shows what it writes instead, and why that is
the same number.

\section{Roots}

\begin{theorem}[Existence and uniqueness of roots]
  For $b > 0$ real and $q \ge 1$ an integer, there is exactly one real $r > 0$ with $r^q = b$.
\end{theorem}
\begin{proof}
  Existence is the intermediate value theorem applied to $t \mapsto t^q$ on $[0, \max(1, b)]$,
  which is continuous, $0$ at $0$ and at least $b$ at the right endpoint. Uniqueness: $t \mapsto
  t^q$ is strictly increasing on $[0, \infty)$, since $0 \le s < t$ gives $s^q < t^q$ by
  induction on $q$.
\end{proof}

Write $b^{1/q}$ for this $r$, and $b^{p/q} = (b^{1/q})^p$. That this agrees with
$\exp(\frac{p}{q}\ln b)$ is the statement that both are the unique positive $q$-th root of $b^p$.
For $b < 0$ and odd $q$ there is a real root, but it is not what $\exp(x \ln |b|)\cos \pi x$
gives (Chapter~\ref{ch:exp}), so the engine restricts its radical rules to integer bases at
least $2$, where every convention agrees.

\section{Square-free parts}

\begin{theorem}[Fundamental theorem of arithmetic]
  Every integer $n \ge 2$ is a product of primes, uniquely up to order.
\end{theorem}

Existence is strong induction: $n$ is prime or a product of two smaller numbers. Uniqueness is
Euclid's lemma again (Chapter~\ref{ch:q}): if a prime $p$ divides $ab$ then $p \mid a$ or
$p \mid b$; so two factorizations of the same $n$ can be matched prime by prime.

\begin{definition}[$q$-th-power part]
  For $b \ge 1$ and $q \ge 2$, write $b = m^q s$ with $m$ as large as possible. Then $s$ is
  \emph{$q$-th-power-free}: no prime occurs in $s$ to a power $\ge q$. For $q = 2$, $s$ is the
  square-free part of $b$.
\end{definition}

By unique factorization, if $b = \prod p_i^{e_i}$ then $m = \prod p_i^{\lfloor e_i/q \rfloor}$
and $s = \prod p_i^{e_i \bmod q}$, and the decomposition is unique. So $8 = 2^2 \cdot 2$ and
$\sqrt 8 = 2\sqrt 2$; $50 = 5^2 \cdot 2$ and $18 = 3^2 \cdot 2$, so $\sqrt{50} - \sqrt{18} =
5\sqrt 2 - 3\sqrt 2 = 2\sqrt 2$ --- the radicals collect because their square-free parts agree.

\begin{theorem}[Irrationality]
  If $b \ge 2$ is not a perfect $q$-th power, then $b^{1/q}$ is irrational.
\end{theorem}
\begin{proof}
  Suppose $b^{1/q} = a/c$ in lowest terms. Then $a^q = b c^q$. Every prime in $a^q$ occurs to a
  power divisible by $q$; every prime in $c^q$ likewise; so every prime in $b$ occurs to a power
  divisible by $q$, and $b$ is a perfect $q$-th power.
\end{proof}

For $b = 2$, $q = 2$ this is the proof that $\sqrt 2$ is irrational, in the form that
generalizes.

\section{The general identity}

Every radical rule the engine has is one identity, read three ways. For $b = m^q s$ and a
rational exponent $x = p/q$, write $p = iq + f$ with $0 \le f < q$. Then
\[
  b^{p/q} = (m^q s)^{p/q} = m^p \, s^{p/q} = m^p\, s^{i}\, s^{f/q}.
\]
Every factor but the last is an integer. So a term $k \cdot b^{p/q}$ is $(k m^p s^i) \cdot s^{f/q}$,
a rational coefficient times a radical of a $q$-th-power-free base with an exponent in $[0, 1)$.
Two such terms with the same $s$, $q$ and $f$ add by adding their coefficients; two radicals with
the same $q$ multiply under one root, $a^{p/q} b^{r/q} = (a^p b^r)^{1/q}$.

\section{What the engine writes}

Here the chapter meets the ordering of Chapter~\ref{ch:wf}. The engine's rewriter terminates
because every rule makes a certain weight decrease, and $8^{1/2} \to 2 \cdot 2^{1/2}$ does not
decrease it: the result contains the input's radical \emph{and} a product node. No retuning of
the weight can fix this, and Chapter~\ref{ch:cannot} gives the reason (a numeral's weight
would have to depend on its value, and then constant folding stops decreasing). What the ordering
does allow is the single-power form: $8^{1/2} \to 2^{3/2}$, one power node for one power node,
ordered by a tier that counts the integers' magnitudes.

\begin{minted}{lean}
def radicalBase : PlainRule :=
  { name := "simp.radical", apply := fun e =>
      match e with
      | .pow (.num a) (.num q) =>
        if a.isInt && a.val.num ≥ 2 && !q.isInt then
          match perfectPower a.val.num.toNat with
          | some (r, k) =>
            let res : Expr := .pow (.num (Q.ofInt r)) (.num (q * Q.ofInt k))
            if M res ≤ M e && numCount res < numCount e then some ⟨res, "…", none, none⟩ else none
          | none => none
        else none
      | _ => none }
\end{minted}

The rule reduces a perfect-power base: $a = r^k$ gives $a^{q} = r^{kq}$. It guards itself with the
decrease its termination proof needs, so the proof is a case split on the guard. Two more rules
do the collecting and the multiplying:

\begin{minted}{lean}
def mergeRadicals (s t : Expr) : Option Expr    -- k₁ b₁^(p₁/q) + k₂ b₂^(p₂/q), same s and f
def mulRadicalPair (s t : Expr) : Option Expr   -- a^(p/q) · b^(r/q) = (a^p b^r)^(1/q)
\end{minted}

\lc{mergeRadicals} computes $m$ and $s$ for each base with \lc{qthPowerPart}, checks that the
$q$-th-power-free parts and the fractional parts agree, and adds the coefficients
$k m^p s^i$. The printer then shows $2^{3/2}$ as $2\sqrt 2$ and $12^{1/2}$ as $2\sqrt 3$: the
display is the textbook's, the term is the ordering's, and the identity of this chapter is the
distance between them.

\section{The theorems}

All three rules are \status{verified}, unconditionally, because their bases are positive integers.

\begin{minted}{lean}
theorem perfectPower_spec {a r k : ℕ} (h : perfectPower a = some (r, k)) : r ^ k = a
theorem qthPowerPart_spec (b q : ℕ) :
    (qthPowerPart b q).1 ^ q * (qthPowerPart b q).2 = b ∧ 1 ≤ (qthPowerPart b q).1
theorem rpow_decompose {b : ℝ} (hb : 0 < b) {m s q : ℕ} (hq : q ≠ 0) (hm : 1 ≤ m) …
theorem radicalBase_soundR (ρ : EnvR) {e : Expr} {res : RuleResult}
    (h : radicalBase.apply e = some res) : evalR ρ res.result = evalR ρ e
theorem collectRadicals_soundR (ρ : EnvR) … : evalR ρ res.result = evalR ρ e
theorem mulRadicals_soundR (ρ : EnvR) … : evalR ρ res.result = evalR ρ e
\end{minted}

The first two are about the searches: a perfect power found is one, a decomposition found is one.
Neither says the search is \emph{complete} --- that it finds the largest $m$ --- because
soundness does not need it; a missed decomposition costs a less simplified answer, never a wrong
one. \lc{rpow_decompose} is the identity of the previous section as a statement about
\lc{Real.rpow}, proved from \lc{Real.rpow_natCast}, \lc{Real.rpow_add} and \lc{Real.rpow_mul},
each of which needs the base positive, which \lc{Q_pos_of_ge_two} supplies. The three rule
theorems assemble it with the permutation lemmas, since the pair the rule merged may have been
anywhere in the list.

\begin{trybox}
\begin{minted}{text}
sqrt(8)
sqrt(50) - sqrt(18)
sqrt(12) * sqrt(3)
4^(2/3)
\end{minted}
  The text output of the first is $2^{3/2}$; the rendered output is $2\sqrt 2$. The last is
  $2^{4/3}$, shown as $2\sqrt[3]{2}$.
\end{trybox}

\section{Exercises}

\begin{exercisebox}[By hand]
  Write $\sqrt{360}$, $\sqrt[3]{250}$ and $5\sqrt{12} - \sqrt{75}$ in the form of the general
  identity, computing $m$, $s$, $i$ and $f$ for each.
\end{exercisebox}

\begin{exercisebox}[Completeness]
  \lc{qthPowerPart} searches $m$ downward from $2^{\lfloor \log_2 b / q \rfloor + 1}$. Show that
  the largest $m$ with $m^q \mid b$ is below that bound, so the search is in fact complete, and
  say why the engine did not need to prove it.
\end{exercisebox}

\begin{exercisebox}[Odd roots of negatives]
  $(-8)^{1/3}$ has a real cube root $-2$. Explain why the engine's rules do not fire on it, and
  what \lc{evalR} gives for it.
\end{exercisebox}
